\documentclass[11pt]{article}
\usepackage[margin=1in]{geometry}
\usepackage{amsmath,amssymb}
\usepackage{graphicx}
\usepackage{booktabs}
\usepackage{hyperref}
\usepackage{natbib}
\usepackage{xcolor}

\title{Direction Instability Predicts Cross-Cell-Line Drug Mechanism Transport in LINCS L1000}
\author{Elliot Tower\\
\texttt{elliot@elliottower.ai}}
\date{July 2026}

\begin{document}
\maketitle

\begin{abstract}
A drug's mechanism of action may or may not generalize across cell types.
We adapt a geometric measure---direction instability, the mean pairwise
angular deviation of perturbation signature vectors---to quantify whether
a drug's transcriptomic effect is conserved across biological contexts.
Applied to 8{,}949 compounds tested in $\geq5$ cell lines from the LINCS
L1000 Connectivity Map (978 landmark genes, 167{,}266 signatures), direction
instability separates mechanism-conserving drugs from context-dependent ones
(population-null $p<0.01$ for 78 drugs). Direction instability anticorrelates
with gene-level Jaccard overlap of top perturbed genes (Spearman
$\rho=-0.79$), confirming that directional change reflects biological
mechanism change. Within the HDAC inhibitor class, pan-isoform inhibitors
transport across cell lines (mean instability 0.62, $n=8$) while
isoform-selective inhibitors do not (mean instability 0.96, $n=5$),
demonstrating that target breadth---not drug class membership---determines
cross-context generalizability.
\end{abstract}

\section{Introduction}

Pharmacological perturbation experiments increasingly rely on profiling drug
effects across panels of cell lines to identify compounds with robust
mechanisms of action \citep{subramanian2017next}. The LINCS L1000 Connectivity
Map measures 978 landmark gene expression changes for ${\sim}20{,}000$
compounds across up to 70 cell lines, producing high-dimensional signature
vectors that characterize each drug's transcriptomic effect in each context
\citep{subramanian2017next}.

A fundamental question arises: when does a drug's effect \emph{transport}
from one cell type to another? Two drugs may share a mechanism-of-action
label yet differ in how consistently they perturb gene expression across
contexts. A pan-HDAC inhibitor like vorinostat affects chromatin remodeling
in every cell, while an HDAC8-selective inhibitor like PCI-34051 only
produces a strong transcriptomic signature where HDAC8 is functionally
relevant.

We formalize this distinction using \textbf{direction instability}: the mean
pairwise angular deviation of a drug's unit-normalized signature vectors
across cell lines. This metric is adapted from the neural bracket norm
\citep{tower2026bracket}, which quantifies how much a learned representation's
direction rotates across evidence levels. Here, the analogy maps evidence
quartiles to cell lines, and neural activation vectors to drug perturbation
signatures.

Direction instability decomposes the perturbation effect into two independent
components. The \emph{direction} captures which genes are perturbed and in
what relative proportions. The \emph{magnitude} captures the overall effect
size. A drug with low direction instability but high magnitude variation has
a conserved qualitative mechanism with context-dependent penetrance---a
biologically meaningful and pharmacologically useful category.

We apply this framework to 8{,}949 drugs with $\geq5$ cell lines in LINCS
L1000. The scorer and validation labels were frozen under a pre-registered
commit SHA before any real data were analyzed.

\section{Methods}

\subsection{Data}

Perturbation signatures were drawn from the LINCS L1000 Phase I dataset
(GEO accession GSE92742), using Level 5 replicate-collapsed z-scores across
978 landmark genes \citep{subramanian2017next}. Metadata (signature
identifiers, perturbation types, cell line annotations) were downloaded from
GEO and filtered to small-molecule perturbagens (\texttt{pert\_type =
trt\_cp}) with $\geq5$ distinct cell lines. Replicate signatures within a
drug--cell-line pair were averaged to yield one consensus signature per
drug per cell line.

Mechanism-of-action (MOA) annotations were obtained from the Broad Drug
Repurposing Hub (version 2020-03-24) \citep{corsello2017drug}, providing
MOA labels for 1{,}782 of the 8{,}949 drugs. These labels were frozen
before analysis (commit SHA \texttt{1dc20a2}).

\subsection{Direction instability}

For a drug tested in $K$ cell lines, let $\mathbf{s}_1, \ldots,
\mathbf{s}_K \in \mathbb{R}^{978}$ denote its consensus signature vectors.
Define the unit-normalized signatures $\hat{\mathbf{s}}_i =
\mathbf{s}_i / \|\mathbf{s}_i\|_2$. Direction instability is:
\begin{equation}
  D = 1 - \frac{2}{K(K-1)} \sum_{i < j}
    \hat{\mathbf{s}}_i^\top \hat{\mathbf{s}}_j
  \label{eq:direction_instability}
\end{equation}
$D \in [0, 2]$. $D = 0$ when all signatures point in the same direction.
$D = 1$ when signatures are orthogonal on average. $D > 1$ indicates
anticorrelation.

\subsection{Complementary metrics}

\textbf{Magnitude CV}: coefficient of variation of signature $L_2$ norms
$\{\|\mathbf{s}_i\|\}_{i=1}^K$, capturing effect-size variation independent
of direction.

\textbf{Fr\'echet variance}: mean squared geodesic deviation from the
Fr\'echet mean on the unit sphere $S^{977}$.

\textbf{Gene-level Jaccard overlap}: for each pair of cell lines, the
Jaccard index of the top-50 up-regulated and top-50 down-regulated genes,
averaged over all pairs and both directions.

\subsection{Population null model}

A drug's direction instability depends on its cell-line count: drugs tested
in more cell lines have more pairwise comparisons and tend toward higher
baseline instability. To control for this confound, we compare each drug
against a reference population of drugs with similar coverage: all drugs
with cell-line count in $[K/2,\; 2K]$, where $K$ is the focal drug's count.
The $p$-value is the fraction of reference drugs with instability $\leq$
the focal drug's instability. This replaces a permutation null, which
is invariant under cell-line label shuffling for pairwise cosine statistics.

\subsection{Pre-registration}

The scoring functions (Equation~\ref{eq:direction_instability} and all
complementary metrics) and the frozen MOA labels were committed before any
real LINCS signatures were analyzed. The commit SHA and frozen label file
are available in the repository.

\section{Results}

\subsection{Landscape of drug mechanism transport}

Among 8{,}949 drugs with $\geq5$ cell lines, mean direction instability
was $0.924 \pm 0.061$ (median 0.940). Most drugs exhibit context-dependent
effects. Only 78 drugs (0.9\%) achieved $p < 0.01$ under the population
null, and 411 (4.6\%) achieved $p < 0.05$, indicating that genuinely
mechanism-conserving drugs are rare.

Cell-line coverage ranged from 5 to 64 (mean 10.2, median 8). The
population null accounts for this variation: a drug's instability is
evaluated against drugs with comparable coverage.

\subsection{Direction instability reflects biological mechanism conservation (H3)}

Direction instability anticorrelated with gene-level Jaccard overlap of
top perturbed genes: Spearman $\rho = -0.79$ ($p < 10^{-300}$). Drugs with
low instability literally perturb the same genes across cell types, while
high-instability drugs engage different gene programs in different contexts.
This confirms that direction instability captures genuine biological
mechanism change, not measurement noise.

\subsection{Direction and magnitude instability are partially dissociable (H4)}

Direction instability and magnitude CV showed weak correlation (Spearman
$\rho = -0.10$, $p = 2.1 \times 10^{-21}$). 7.8\% of drugs fell in
the low-direction, high-magnitude quadrant: conserved mechanism with
context-dependent penetrance.

\subsection{MOA class stratification (H2)}

Table~\ref{tab:moa} shows mean direction instability by MOA class (classes
with $\geq10$ annotated drugs). Two classes stand out: HDAC inhibitors
($\bar{D} = 0.77$, $\sigma = 0.14$) and topoisomerase inhibitors
($\bar{D} = 0.77$, $\sigma = 0.12$). CDK inhibitors and tubulin
polymerization inhibitors also show reduced instability ($\bar{D} = 0.80$
and $0.88$ respectively). Receptor antagonists and agonists cluster near
the population mean ($\bar{D} = 0.93$--$0.96$), consistent with
receptor-mediated mechanisms being cell-type-dependent.

\begin{table}[t]
\centering
\caption{Direction instability by MOA class ($n \geq 10$). Classes targeting
fundamental cellular machinery (chromatin, DNA, cell cycle) show lower
instability than receptor-mediated classes. Within HDAC inhibitors, the
high variance ($\sigma = 0.14$) reflects the pan-selective gradient
discussed in Section~\ref{sec:hdac}.}
\label{tab:moa}
\small
\begin{tabular}{lrrr}
\toprule
MOA class & $n$ & $\bar{D}$ & $\sigma$ \\
\midrule
topoisomerase inhibitor & 17 & 0.765 & 0.123 \\
HDAC inhibitor & 20 & 0.766 & 0.142 \\
CDK inhibitor & 16 & 0.797 & 0.119 \\
tubulin polymerization inhibitor & 19 & 0.882 & 0.048 \\
MEK inhibitor & 14 & 0.888 & 0.066 \\
EGFR inhibitor & 23 & 0.899 & 0.064 \\
FLT3 inhibitor & 15 & 0.896 & 0.057 \\
tyrosine kinase inhibitor & 14 & 0.901 & 0.065 \\
VEGFR inhibitor & 21 & 0.909 & 0.045 \\
calcium channel blocker & 25 & 0.923 & 0.067 \\
monoamine oxidase inhibitor & 17 & 0.924 & 0.071 \\
glucocorticoid receptor agonist & 31 & 0.933 & 0.039 \\
\midrule
\emph{Population mean} & 8{,}949 & 0.924 & 0.061 \\
\midrule
dopamine receptor antagonist & 59 & 0.936 & 0.046 \\
histamine receptor antagonist & 41 & 0.944 & 0.049 \\
serotonin receptor antagonist & 62 & 0.945 & 0.035 \\
cyclooxygenase inhibitor & 53 & 0.947 & 0.031 \\
adrenergic receptor antagonist & 52 & 0.945 & 0.038 \\
adrenergic receptor agonist & 40 & 0.950 & 0.028 \\
phosphodiesterase inhibitor & 31 & 0.953 & 0.023 \\
serotonin receptor agonist & 31 & 0.958 & 0.022 \\
PPAR receptor agonist & 19 & 0.956 & 0.019 \\
\bottomrule
\end{tabular}
\end{table}

\subsection{The HDAC inhibitor selectivity gradient}
\label{sec:hdac}

The HDAC inhibitor class provides a natural experiment: 20 annotated HDAC
inhibitors span the full selectivity spectrum from pan-isoform to
isoform-specific. Table~\ref{tab:hdac} shows the eight most-transporting
HDAC inhibitors (all pan-isoform) and the five least-transporting
(selective or weak).

\begin{table}[t]
\centering
\caption{HDAC inhibitors ranked by direction instability. Pan-isoform
inhibitors achieve low instability and significant population $p$-values;
selective inhibitors are indistinguishable from the null.}
\label{tab:hdac}
\small
\begin{tabular}{llrrr}
\toprule
Drug & Selectivity & $D$ & $n$ & $p$ \\
\midrule
PCI-24781 & pan & 0.561 & 13 & 0.004 \\
belinostat & pan & 0.589 & 13 & 0.006 \\
panobinostat & pan & 0.591 & 15 & 0.007 \\
scriptaid & pan & 0.620 & 15 & 0.009 \\
vorinostat & pan & 0.628 & 60 & 0.008 \\
trichostatin-a & pan & 0.634 & 64 & 0.014 \\
NSC-3852 & pan & 0.674 & 15 & 0.016 \\
entinostat & class I & 0.697 & 14 & 0.023 \\
\midrule
valproic acid & weak/partial & 0.955 & 54 & 0.947 \\
Merck60 & selective & 0.962 & 9 & 0.778 \\
PCI-34051 & HDAC8-selective & 0.984 & 12 & 0.976 \\
phenylbutyrate & weak/partial & 0.989 & 11 & 0.988 \\
\bottomrule
\end{tabular}
\end{table}

Pan-HDAC inhibitors (mean $D = 0.62$, $n = 8$) transport their mechanism
across cell lines because their targets---histone deacetylases of classes
I, II, and IV---are expressed in every cell type. Chromatin remodeling is a
universal process, so inhibiting it broadly produces a conserved
transcriptomic signature.

Selective HDAC inhibitors (mean $D = 0.96$, $n = 5$) do not transport.
PCI-34051 targets HDAC8 specifically; its transcriptomic effect depends on
whether HDAC8 is functionally relevant in a given cell type. Valproic acid
and phenylbutyrate are weak HDAC inhibitors with additional pharmacological
activities, so their signatures reflect cell-type-specific off-target
effects rather than HDAC inhibition per se.

Entinostat ($D = 0.70$, class I--selective) falls between the two groups,
consistent with intermediate target breadth: class I HDACs (HDACs 1, 2, 3)
are more widely expressed than HDAC8 but narrower than the full isoform
panel.

\subsection{Top transporting drugs target fundamental machinery}

The 10 most mechanism-conserving drugs with MOA annotations target processes
common to all cells (Table~\ref{tab:top}): protein synthesis
(homoharringtonine), RNA transcription (dactinomycin), DNA topology
(epirubicin, doxorubicin), chromatin remodeling (PCI-24781, belinostat,
panobinostat), cell cycle progression (BMS-387032), and signal transduction
kinases (bisindolylmaleimide-IX). These mechanisms operate independently
of tissue-specific receptor expression, explaining their cross-context
conservation.

\begin{table}[t]
\centering
\caption{Top 10 transporting drugs with MOA annotation.}
\label{tab:top}
\small
\begin{tabular}{llrr}
\toprule
Drug & MOA & $D$ & $p$ \\
\midrule
homoharringtonine & protein synthesis inhibitor & 0.497 & 0.001 \\
dactinomycin & RNA polymerase inhibitor & 0.506 & 0.002 \\
BMS-387032 & CDK / cell cycle inhibitor & 0.522 & 0.002 \\
bisindolylmaleimide-IX & PKC inhibitor & 0.526 & 0.003 \\
epirubicin & topoisomerase inhibitor & 0.529 & 0.003 \\
anisomycin & DNA synthesis inhibitor & 0.555 & 0.004 \\
PCI-24781 & HDAC inhibitor & 0.561 & 0.004 \\
BNTX & opioid receptor antagonist & 0.564 & 0.003 \\
digitoxigenin & ATPase inhibitor & 0.586 & 0.004 \\
belinostat & HDAC inhibitor & 0.589 & 0.006 \\
\bottomrule
\end{tabular}
\end{table}

Conversely, imatinib ($D = 0.97$) exemplifies a context-dependent drug:
its BCR-ABL kinase inhibition produces strong effects only in cells
harboring the Philadelphia chromosome, yielding high direction instability
across the LINCS panel.

\section{Discussion}

Direction instability provides a single-number summary of whether a drug's
transcriptomic mechanism generalizes across cell types. Three findings
support its validity.

First, the metric tracks gene-level mechanism conservation ($\rho = -0.79$
with Jaccard overlap). This is a strong sanity check: if direction
instability were driven by noise or batch effects, it would not correlate
with which specific genes are perturbed.

Second, MOA classes stratify in a biologically interpretable pattern.
Drugs targeting universal cellular machinery (chromatin remodeling, DNA
topology, protein synthesis, cell cycle) show low instability, while
receptor-mediated drugs show high instability. This matches the
expectation that receptor expression is tissue-specific.

Third, the HDAC inhibitor selectivity gradient provides a within-class
controlled comparison. Pan-isoform and selective HDAC inhibitors share
an MOA label but differ in target breadth, and this difference maps
directly onto direction instability. This rules out the alternative
explanation that MOA labels alone drive the signal.

\paragraph{Relation to the bracket norm.}
Direction instability is formally analogous to the neural bracket norm
\citep{tower2026bracket}, which measures directional rotation of
learned representations across evidence levels. Both metrics detect when a
system (neural network or cell) responds to a perturbation (evidence or
drug) by rotating its internal state rather than scaling it. The
pharmacological application replaces evidence quartiles with cell lines
and neural activations with gene expression signatures, preserving the
geometric interpretation: low rotation implies a conserved mechanism that
transports across contexts.

\paragraph{Limitations.}
LINCS L1000 measures only 978 landmark genes; full-transcriptome profiling
might reveal additional structure. The population null controls for
cell-line count but not for other confounds (drug concentration, time point,
signature quality). MOA annotations from the Repurposing Hub are incomplete
(20\% coverage) and may contain errors. The direction-instability metric
treats all cell lines as exchangeable and does not model phylogenetic
relationships between cell types.

\paragraph{Applications.}
Direction instability could serve as a pre-clinical filter for drug
repurposing: compounds with low instability are more likely to produce
consistent effects across patient populations. The metric requires only
perturbation signatures across $\geq5$ contexts, making it applicable to
any multi-context profiling dataset (CMAP, JUMP-CP, Perturb-seq).

\bibliographystyle{plainnat}
\bibliography{references}

\end{document}
